\documentclass[conference]{IEEEtran}
\IEEEoverridecommandlockouts
\usepackage{cite}
\usepackage{amsmath,amssymb,amsfonts}
\usepackage{graphicx}
\usepackage{textcomp}
\usepackage{xcolor}
\usepackage{booktabs}
\usepackage{array}
\usepackage{microtype}
\usepackage[hidelinks]{hyperref}

\setlength{\emergencystretch}{2em}

\begin{document}

\title{Lightweight AI-Based Small Drone Detection, Visual Tracking, and Short-Term Trajectory Prediction Under Occlusion and Camera Motion}

\author{
\IEEEauthorblockN{Javeria Dawood}
\IEEEauthorblockA{\textit{Dept. of Electrical Engineering} \\
\textit{Bahria School of Eng. \& Applied Sciences}\\
Bahria University, Islamabad, Pakistan \\
javeriadawood433@gmail.com}
\and
\IEEEauthorblockN{Syed Ali Turab Kazmi}
\IEEEauthorblockA{\textit{Dept. of Electrical Engineering} \\
\textit{Bahria School of Eng. \& Applied Sciences}\\
Bahria University, Islamabad, Pakistan \\
syedaliturab072@gmail.com}
\and
\IEEEauthorblockN{Amber Saeed}
\IEEEauthorblockA{\textit{Dept. of Electrical Engineering} \\
\textit{Bahria School of Eng. \& Applied Sciences}\\
Bahria University, Islamabad, Pakistan \\
ambersaeed2005@gmail.com}
}

\maketitle

\begin{abstract}
Low-altitude micro-unmanned aerial vehicles (UAVs) are difficult to monitor reliably with onboard or ground-based cameras when computational power is limited. In real-world flights, four problems occur together: distant quadrotors subtend fewer than $24 \times 24$ pixels, urban buildings and tree foliage lower detector confidence, obstacles briefly block the line of sight, and camera pan or tilt motion shifts bounding boxes across the image plane. In this paper, we develop and test a lightweight aerial tracking pipeline that couples a modified YOLOv8n-P2 small-target detector, two-stage ByteTrack data association, monocular Visual-Inertial Odometry (VIO), and a 3D Constant-Turn-Rate and Acceleration Extended Kalman Filter (CTRA-EKF). A high-resolution $160 \times 160$ P2 feature head with Wise-IoU regression captures sub-20-pixel drones, while tightly coupled VIO fuses $30\text{ Hz}$ visual features with $200\text{ Hz}$ inertial measurement unit (IMU) pre-integration to compensate camera ego-rotation and translation prior to ByteTrack association. During complete visual occlusions ($1.0\text{--}3.0\text{ s}$), the 3D CTRA-EKF propagates curvilinear target dynamics to prevent track fragmentation. Evaluated across 16 physics-based aerial surveillance sequences ($14,400$ frames) in Gazebo Harmonic 8.15 alongside standard UAV small-object benchmarks, our framework achieves $91.4\%$ $\text{AP}_{50}$, $89.4\%$ MOTA, $91.2\%$ IDF1, and $81.6\%$ HOTA while reducing identity switches by $78.9\%$ over uncompensated baselines. Camera ego-motion estimation reaches an Absolute Trajectory Error ($\text{ATE}_{\text{RMSE}}$) of $0.084\text{ m}$, while the 3D CTRA-EKF lowers $2.0\text{ s}$ blind-occlusion prediction RMSE by $81.8\%$ ($0.36\text{ m}$ vs. $1.98\text{ m}$). The full pipeline runs at $128.2\text{ FPS}$ ($7.8\text{ ms}$ per frame) on a desktop GPU and $46.1\text{ FPS}$ on an embedded Jetson Orin Nano. All simulation worlds, ROS~2/Gazebo nodes, and evaluation datasets are publicly available at \url{https://github.com/T-Storm-072/03_uav_drone_tracking_icect2027}.
\end{abstract}

\begin{IEEEkeywords}
Small drone detection, YOLOv8, ByteTrack, Extended Kalman Filter (EKF), Visual-Inertial Odometry (VIO), occlusion handling, Gazebo simulation.
\end{IEEEkeywords}

\section{Introduction}
Small rotary-wing Unmanned Aerial Vehicles (UAVs) are now routine in crop monitoring, bridge inspection, search-and-rescue, and site security \cite{ref1, ref2}. At the same time, low-cost micro-quadrotors have made airspace monitoring much harder, requiring vision systems that can spot, track, and predict fast aerial trajectories on low-power hardware \cite{ref3}.

When a camera tracks a small drone outdoors, four failure modes interact directly. First, beyond $25\text{ m}$ standoff range, a $0.28\text{ m}$ quadrotor covers only $8 \times 8$ to $24 \times 24$ pixels ($<0.08\%$ of a $1280 \times 720$ image), so standard stride-8 and stride-16 backbone layers lose its spatial features \cite{ref1, ref4}. Second, brick facades, moving branches, and changing sky brightness cause confidence scores to drop and trigger false alarms \cite{ref5, ref6}. Third, when a drone passes behind a pillar or tree canopy, visual measurements disappear for several seconds; standard constant-velocity Kalman filters then drift along the pre-occlusion tangent and assign a new track ID when the drone reappears \cite{ref9, ref10}. Fourth, if the camera itself rotates on a gimbal or flies onboard a pursuit UAV, camera pitch and yaw shifts move bounding boxes across the frame even when the target hovers still \cite{ref11, ref12}.

While recent studies have independently advanced lightweight YOLO architectures for small aerial targets \cite{ref1, ref3, ref7, ref8}, multi-stage detection association via ByteTrack \cite{ref9}, and monocular Visual-Inertial Odometry (VIO) via VINS-Mono \cite{ref12}, existing literature largely evaluates these modules in isolation. Static-camera or downward-looking ground-target benchmarks (such as standard VisDrone pedestrian/vehicle tasks \cite{ref2}) do not establish how detector confidence drops, camera ego-motion drift, and multi-second blind occlusions interact during dynamic air-to-air or ground-to-air drone tracking.

In this work, we address these four failure modes together inside a unified ROS 2 and Gazebo Harmonic 8.15 simulation framework. Specifically, our contributions are:
\begin{itemize}
    \item \textbf{High-Resolution Small-UAV Detector:} An enhanced YOLOv8n-P2 architecture incorporating a $160 \times 160$ shallow feature pyramid branch and dynamic Wise-IoU bounding-box regression tailored for sub-24-pixel drones.
    \item \textbf{VIO Ego-Motion Compensated ByteTrack:} Tightly coupled integration of $200\text{ Hz}$ IMU pre-integration and monocular visual odometry ($\mathbf{T}_{w}^{c_t} \in SE(3)$) to project tracklet priors across aggressive camera pan/tilt jitter prior to two-stage ByteTrack association.
    \item \textbf{Curvilinear 3D CTRA-EKF Occlusion Predictor:} A nonlinear Extended Kalman Filter modeling 3D target velocity, yaw angle, and turn rate to bridge $1.0\text{--}3.0\text{ s}$ complete visual occlusions with minimal spatial drift.
    \item \textbf{Comprehensive Gazebo \& Benchmark Evaluation:} Rigorous ablation across 16 simulated sequences ($14,400$ frames) and 4 environmental regimes ($R_1\text{--}R_4$), quantifying accuracy ($\text{AP}_{50}$, MOTA, IDF1, HOTA, $\text{RMSE}$, $\text{ATE}_{\text{RMSE}}$) and real-time computational efficiency (FPS, latency, and memory).
\end{itemize}

\begin{figure*}[!t]
\centering
\includegraphics[width=0.88\textwidth]{fig1_uav_arch.png}\vspace{-6pt}
\caption{Proposed integrated aerial monitoring architecture combining YOLOv8n-P2 small-drone detection, tightly coupled VINS-Mono Visual-Inertial Odometry (VIO) ego-motion compensation, two-stage ByteTrack association, and a 3D CTRA Extended Kalman Filter (EKF) for short-term trajectory prediction under occlusion.}
\label{fig:arch}
\end{figure*}

\section{Literature Review}

\subsection{Lightweight Small-Object Detection in Aerial Imagery}
Objects occupying only a few dozen pixels suffer from spatial information loss during strided downsampling in deep convolutional networks \cite{ref2}. Xu et al. \cite{ref1} improved YOLOv8n on the VisDrone benchmark by introducing an additional small-scale detection layer, conditionally parameterized convolutions, and a Wise-IoU loss function, reaching $37.6\%$ mAP with $3.0\text{ M}$ parameters at $40\text{ FPS}$. Li et al. \cite{ref3} proposed SOD-YOLO, integrating a Receptive Field Convolutional Block Attention Module (RFCBAM) and Balanced Spatial-Semantic Information Fusion Pyramid Network (BSSI-FPN) to boost small-target $\text{AP}_{50}$ while cutting parameter counts. Similarly, Xu et al. \cite{ref4} developed YOLOv8-MPEB by replacing CSPDarknet53 with a MobileNetV3 backbone ($7.39\text{ M}$ parameters), while Feng et al. \cite{ref5} integrated Coordinate Attention and a Dynamic Head to reduce missed detections in cluttered backgrounds.

More recently, Xiao and Di \cite{ref6} introduced a Deep-Shallow Downsampling Module (DSDM) with a P2 detection head in Scientific Reports, achieving $50.7\%$ $\text{AP}_{50}$ at $72.5\text{ FPS}$. Liang et al. \cite{ref7} proposed Lightweight YOLO-SR using Shuffle Attention (`C2f-SA`) and decoupled heads (`DDHead`), reaching $86.5\%$ $\text{AP}_{50}$ on RGBT-Tiny. Jin et al. \cite{ref8} presented GS-YOLO in PLOS ONE, combining `C2FGhostLight` and Small-Object-Aware Lightweight Channel Attention (`SOLCA`) to reduce parameters to $0.84\text{ M}$. While these works confirm the value of P2 shallow heads and attention-guided fusion, they primarily evaluate static-frame ground targets rather than continuous 3D drone trajectory tracking under camera motion.

\subsection{Multi-Object Visual Tracking Under Occlusion}
Detection alone yields unlinked per-frame bounding boxes; visual tracking must preserve target identity across time despite confidence fluctuations and partial occlusion \cite{ref9, ref10}. Zhang et al. \cite{ref9} introduced ByteTrack, which resolves the fragmentation caused by discarding low-score detections. By matching high-confidence detections ($s \ge \tau_{\text{high}}$) in a first pass and recovering low-score boxes ($\tau_{\text{low}} \le s < \tau_{\text{high}}$) against remaining tracklets in a second pass, ByteTrack raises both IDF1 \cite{ref17} and HOTA \cite{ref16} on multi-object tracking benchmarks. Cao et al. \cite{ref10} further showed in OC-SORT that linear constant-velocity Kalman filters accumulate large direction errors during occlusion gaps, motivating observation-centric momentum and nonlinear motion models, while Aharon et al. \cite{ref11} demonstrated that uncompensated camera rotation degrades IoU matching.

\subsection{Camera-Motion Estimation and Visual-Inertial Systems}
When the observing camera moves, target displacement in the image plane is the superposition of true 3D target velocity and camera ego-motion. Qin et al. \cite{ref12} formulated VINS-Mono, a tightly coupled nonlinear optimization framework fusing monocular visual features with pre-integrated IMU measurements to estimate 6-DoF camera pose $\mathbf{T}_{w}^{c_t} \in SE(3)$ in real time. While VINS-Mono accurately tracks camera ego-motion (evaluated via Absolute Trajectory Error $\text{ATE}$ and Relative Pose Error $\text{RPE}$ \cite{ref18}), it does not track independent dynamic targets. Coupling VIO pose estimates directly with ByteTrack and a 3D nonlinear EKF \cite{ref13, ref14} allows the tracker to separate camera rotation from actual 3D target maneuvers.

\section{Proposed Integrated Methodology}
Fig.~\ref{fig:arch} illustrates the proposed four-stage pipeline. At each frame $t_k$, synchronized RGB imagery ($\mathbf{I}_k$ at $30\text{ Hz}$) and high-rate IMU measurements ($\boldsymbol{\omega}_m, \mathbf{a}_m$ at $200\text{ Hz}$) are processed sequentially.

\subsection{YOLOv8n-P2 Small-Drone Detection}
Standard YOLOv8n utilizes feature pyramid outputs at strides $\{8, 16, 32\}$ ($P_3, P_4, P_5$), where a $12 \times 12\text{ px}$ drone is compressed to $1.5 \times 1.5$ cells in $P_3$. Following \cite{ref1, ref6}, we incorporate a stride-4 $P_2$ detection head ($160 \times 160$ resolution for $640 \times 640$ input) fused via bidirectional PAN-FPN pathways with lightweight C2f-Ghost blocks \cite{ref8} and prune redundant deep channels to keep the total parameter count at $2.84\text{ M}$. Bounding-box regression is optimized using Wise-IoU (WIoU) loss \cite{ref1}:
\begin{equation}
\mathcal{L}_{\text{WIoU}} = \mathcal{R}_{\text{WIoU}} \cdot \mathcal{L}_{\text{IoU}}, \quad \mathcal{R}_{\text{WIoU}} = \exp\left(\frac{(x - x_{\text{gt}})^2 + (y - y_{\text{gt}})^2}{(W_g^2 + H_g^2)^*}\right)
\label{eq:wiou}
\end{equation}
which dynamically attenuates harmful gradients from motion-blurred small-drone annotations.

\subsection{Tightly Coupled VIO \& Ego-Motion Compensation}
Let the camera platform state in the world frame $\mathcal{W}$ be $\boldsymbol{\chi}_c = [\mathbf{p}_w^c, \mathbf{v}_w^c, \mathbf{q}_w^c, \mathbf{b}_a, \mathbf{b}_g]^\top$. Between consecutive video frames $[t_{k-1}, t_k]$, $200\text{ Hz}$ IMU measurements are pre-integrated following Qin et al. \cite{ref12} and optimized over a sliding window with corner features excluded from detected drone bounding boxes. Given the resulting inter-frame camera motion $\mathbf{T}_{c_k}^{c_{k-1}} = (\mathbf{T}_w^{c_k})^{-1}\mathbf{T}_w^{c_{k-1}} = [\mathbf{R}_{k-1}^k \mid \mathbf{t}_{k-1}^k]$ and intrinsic matrix $\mathbf{K}$, predicted tracklet center coordinates $\tilde{\mathbf{u}}_{k-1} = [u, v, 1]^\top$ at estimated depth $\hat{d}_{k-1}$ are warped into the current camera frame prior to data association:
\begin{equation}
\hat{\mathbf{u}}_{k|k-1}^{\text{ego}} \propto \mathbf{K}\left(\mathbf{R}_{k-1}^k \hat{d}_{k-1}\mathbf{K}^{-1}\tilde{\mathbf{u}}_{k-1} + \mathbf{t}_{k-1}^k\right)
\label{eq:ego_warp}
\end{equation}
eliminating apparent bounding-box jumps induced by camera pitch, yaw, or translation.

\subsection{Confidence-Stratified ByteTrack Association}
Detections $\mathcal{D}_k$ from YOLOv8n-P2 are partitioned into high-confidence boxes $\mathcal{D}_{\text{high}} = \{b_j \mid s_j \ge \tau_{\text{high}}\}$ ($\tau_{\text{high}} = 0.50$) and low-confidence boxes $\mathcal{D}_{\text{low}} = \{b_j \mid \tau_{\text{low}} \le s_j < \tau_{\text{high}}\}$ ($\tau_{\text{low}} = 0.15$) \cite{ref9}. In Stage~1, active tracklets $\mathcal{T}$ (warped via Eq.~\ref{eq:ego_warp}) are matched against $\mathcal{D}_{\text{high}}$ via the Hungarian algorithm using a hybrid cost matrix combining Generalized IoU and normalized center distance. In Stage~2, remaining unmatched tracklets $\mathcal{T}_{\text{rem}}$ (typically drones partially obscured by foliage or motion blur) are matched against $\mathcal{D}_{\text{low}}$, preserving track continuity before handing off unmatched active tracks to the EKF occlusion predictor.

\begin{figure*}[!t]
\centering
\includegraphics[width=0.84\textwidth]{fig2_uav_gazebo.png}\vspace{-8pt}
\caption{Synchronous aerial simulation keyframes in Gazebo Harmonic 8.15 across all four experimental regimes: (a) $R_1$: Long-range sub-20-pixel micro-UAV detection via the P2 head, (b) $R_2$: ByteTrack Stage-2 low-confidence recovery ($s=0.24$) across foliage clutter, (c) $R_3$: 3D CTRA-EKF short-term trajectory prediction bridging a $2.2\text{ s}$ full building occlusion vs. linear KF tangential drift, and (d) $R_4$: VIO ego-motion flow compensation under $25^\circ/\text{s}$ camera pitch.}
\label{fig:gazebo}
\end{figure*}

\subsection{3D CTRA Extended Kalman Filter (EKF) for Short-Term Trajectory Prediction Under Occlusion}
Unlike terrestrial pedestrians, agile quadrotors execute banked curvilinear turns. We formulate a 3D Constant-Turn-Rate and Acceleration (CTRA) state vector in the world frame $\mathcal{W}$ \cite{ref13, ref14}:
\begin{equation}
\mathbf{x}_k = [x_k, y_k, z_k, v_{xy,k}, \psi_k, \omega_{\psi,k}, a_{xy,k}, v_{z,k}]^\top \in \mathbb{R}^8
\label{eq:state}
\end{equation}
where $(x_k, y_k, z_k)$ is the 3D UAV position, $v_{xy,k}$ is horizontal speed, $\psi_k$ is heading yaw, $\omega_{\psi,k}$ is yaw rate, $a_{xy,k}$ is tangential acceleration, and $v_{z,k}$ is vertical climb rate. The nonlinear discrete-time transition model $\mathbf{x}_{k|k-1} = f(\hat{\mathbf{x}}_{k-1|k-1}, \Delta t)$ for $\omega_{\psi} \neq 0$ is:
\begin{align}
x_{k|k-1} &= x_{k-1} + \frac{v_{xy} + a_{xy}\Delta t}{\omega_\psi}\sin(\psi + \omega_\psi\Delta t) - \frac{v_{xy}}{\omega_\psi}\sin\psi \nonumber\\
&\quad + \frac{a_{xy}}{\omega_\psi^2}\left[\cos(\psi + \omega_\psi\Delta t) - \cos\psi\right] \label{eq:ekf_pred}\\
y_{k|k-1} &= y_{k-1} - \frac{v_{xy} + a_{xy}\Delta t}{\omega_\psi}\cos(\psi + \omega_\psi\Delta t) + \frac{v_{xy}}{\omega_\psi}\cos\psi \nonumber\\
&\quad + \frac{a_{xy}}{\omega_\psi^2}\left[\sin(\psi + \omega_\psi\Delta t) - \sin\psi\right]
\end{align}
with $z_{k|k-1} = z_{k-1} + v_{z,k-1}\Delta t$. State covariance is propagated via the linearized Jacobian $\mathbf{F}_k = \left.\frac{\partial f}{\partial \mathbf{x}}\right|_{\hat{\mathbf{x}}_{k-1|k-1}}$:
\begin{equation}
\mathbf{P}_{k|k-1} = \mathbf{F}_k \mathbf{P}_{k-1|k-1} \mathbf{F}_k^\top + \mathbf{Q}_k
\label{eq:ekf_cov}
\end{equation}
When a target is visible, back-projected bounding-box ray and scale observations $\mathbf{z}_k = h(\mathbf{x}_k) + \boldsymbol{\nu}_k$ ($\boldsymbol{\nu}_k \sim \mathcal{N}(\mathbf{0}, \mathbf{R}_k)$) update the state via Kalman gain $\mathbf{K}_k = \mathbf{P}_{k|k-1}\mathbf{H}_k^\top(\mathbf{H}_k\mathbf{P}_{k|k-1}\mathbf{H}_k^\top + \mathbf{R}_k)^{-1}$. When total occlusion occurs ($\mathcal{D}_{\text{high}} \cup \mathcal{D}_{\text{low}} = \emptyset$ for tracklet $i$), the measurement update is gated off and Eqs.~(\ref{eq:ekf_pred})--(\ref{eq:ekf_cov}) recursively predict the drone's curvilinear path for up to $\tau_{\text{max}} = 3.0\text{ s}$ ($90$ frames), allowing immediate re-identification when the target emerges from behind obstacles.

\section{Simulation Setup \& Evaluation Metrics}

\subsection{Gazebo Harmonic 8.15 Aerial Environment}
Experiments were conducted in Gazebo Harmonic 8.15 coupled with ROS~2 Humble and DART physics (Fig.~\ref{fig:gazebo}). The testbed comprises a moving observer platform ($30\text{ Hz}$ RGB camera, $1280 \times 720$ resolution, $78^\circ$ HFOV; $200\text{ Hz}$ 6-axis IMU with gyroscope noise $\sigma_g = 1.5 \times 10^{-3}\text{ rad/s/}\sqrt{\text{Hz}}$ and accelerometer noise $\sigma_a = 1.8 \times 10^{-2}\text{ m/s}^2/\sqrt{\text{Hz}}$) monitoring agile micro-UAVs ($0.28\text{ m}$ wingspan) across 16 sequences ($14,400$ annotated frames) divided into four regimes:
\begin{itemize}
    \item $R_1$ \textbf{(Long-Range Small Scale):} Target distances of $25\text{--}60\text{ m}$ ($8 \times 8$ to $22 \times 22\text{ px}$).
    \item $R_2$ \textbf{(Complex Urban/Foliage Background):} Flights across textured buildings and tree canopies.
    \item $R_3$ \textbf{(Dynamic Occlusion):} Structured blind intervals of $1.0\text{ s}$, $2.0\text{ s}$, and $3.0\text{ s}$ behind pillars and foliage.
    \item $R_4$ \textbf{(Aggressive Camera Ego-Motion):} Observer rotational rates up to $28^\circ/\text{s}$ and translational acceleration up to $2.5\text{ m/s}^2$.
\end{itemize}

\begin{table*}[!t]
\centering
\caption{Small-Drone Detection Performance Comparison Across Lightweight YOLO Architectures ($14,400$ Validation Frames, Input $640 \times 640$).}
\label{tab:detection}
\footnotesize
\renewcommand{\arraystretch}{1.15}
\resizebox{\textwidth}{!}{\begin{tabular}{@{}lcccccccc@{}}
\toprule
\textbf{Detection Model} & \textbf{Backbone / Head} & \textbf{Params (M)} & \textbf{GFLOPs} & \textbf{Precision (\%)} & \textbf{Recall (\%)} & \textbf{F1-Score (\%)} & \textbf{$\text{AP}_{50}$ (\%)} & \textbf{$\text{mAP}_{50:95}$ (\%)} \\
\midrule
Baseline YOLOv8n \cite{ref15} & CSPDarknet (P3--P5) & 3.01 & 8.1 & 81.4 & 73.2 & 77.1 & 78.6 & 46.2 \\
YOLOv8-MPEB \cite{ref4} & MobileNetV3 (P3--P5) & 7.39 & 11.4 & 84.9 & 77.5 & 81.0 & 82.9 & 49.8 \\
SOD-YOLO \cite{ref3} & RFCBAM + BSSI-FPN & 2.65 & 9.8 & 88.2 & 83.6 & 85.8 & 87.8 & 54.1 \\
Lightweight YOLO-SR \cite{ref7} & C2f-SA + DDHead (P1--P3) & 4.02 & 12.1 & 88.7 & 84.1 & 86.3 & 88.2 & 55.0 \\
GS-YOLO \cite{ref8} & C2FGhostLight + SOLCA & \textbf{0.84} & \textbf{4.6} & 86.5 & 80.9 & 83.6 & 85.1 & 51.9 \\
\textbf{Proposed YOLOv8n-P2} & \textbf{C2f-Ghost + P2 Head + WIoU} & 2.84 & 9.2 & \textbf{92.3} & \textbf{88.7} & \textbf{90.5} & \textbf{91.4} & \textbf{58.9} \\
\bottomrule
\end{tabular}}
\end{table*}

\subsection{Evaluation Metrics}
Following established small-object detection, tracking, and visual-inertial standards \cite{ref6, ref16, ref17, ref18, ref19}, system performance is evaluated across four metric groups:

\subsubsection{Detection Accuracy and Speed}
Precision ($P$), Recall ($R$), harmonic $\text{F1-Score}$, Average Precision ($\text{AP}$), Mean Average Precision ($\text{mAP}$), and Frames Per Second ($\text{FPS}$) are computed from True Positives ($\text{TP}$), False Positives ($\text{FP}$), and False Negatives ($\text{FN}$):
\begin{equation}
P = \frac{\text{TP}}{\text{TP} + \text{FP}}, \quad R = \frac{\text{TP}}{\text{TP} + \text{FN}}, \quad \text{F1} = \frac{2 P R}{P + R}
\label{eq:prf1}
\end{equation}
\begin{equation}
\text{AP} = \int_{0}^{1} P(r)\,dr, \quad \text{mAP} = \frac{1}{C}\sum_{i=1}^{C}\text{AP}_i, \quad \text{FPS} = \frac{N_{\text{frames}}}{T_{\text{total}}}
\label{eq:map_fps}
\end{equation}

\subsubsection{Multi-Object Tracking Consistency}
Identity preservation is quantified via Multi-Object Tracking Accuracy ($\text{MOTA}$) \cite{ref19}, Identity F1 ($\text{IDF1}$) \cite{ref17}, Higher Order Tracking Accuracy ($\text{HOTA}$) \cite{ref16}, and Identity Switches ($\text{IDSW}$):
\begin{equation}
\text{MOTA} = 1 - \frac{\sum_{t}(\text{FN}_t + \text{FP}_t + \text{IDSW}_t)}{\sum_{t}\text{GT}_t}
\label{eq:mota}
\end{equation}

\subsubsection{EKF State Estimation Error}
Short-term trajectory prediction accuracy during occlusion windows of $N$ steps is evaluated via 3D Root Mean Square Error ($\text{RMSE}$):
\begin{equation}
\text{RMSE}_{\text{EKF}} = \sqrt{\frac{1}{N}\sum_{k=1}^{N}\left\|\hat{\mathbf{p}}_{\text{uav},k} - \mathbf{p}_{\text{uav},k}^{\text{gt}}\right\|_2^2}
\label{eq:rmse}
\end{equation}

\subsubsection{VIO Camera Ego-Motion Error}
Global and local ego-motion estimation accuracy are measured using translational Absolute Trajectory Error ($\text{ATE}_{\text{RMSE}}$) and Relative Pose Error ($\text{RPE}$) after $SE(3)$ alignment \cite{ref18}:
\begin{equation}
\text{ATE}_{\text{RMSE}} = \sqrt{\frac{1}{N}\sum_{i=1}^{N}\left\|\mathbf{p}_{c,i}^{\text{gt}} - \hat{\mathbf{p}}_{c,i}^{\text{vio}}\right\|_2^2}
\label{eq:ate}
\end{equation}

\begin{table*}[!t]
\centering
\caption{Ablation Study of the Integrated Detection, Tracking, VIO Ego-Motion Compensation, and 3D CTRA-EKF Framework Across Regimes $R_1\text{--}R_4$.}
\label{tab:ablation}
\footnotesize
\renewcommand{\arraystretch}{1.15}
\resizebox{\textwidth}{!}{\begin{tabular}{@{}llccccccccc@{}}
\toprule
\textbf{Config} & \textbf{Integrated Pipeline Modules} & \textbf{MOTA (\%)} & \textbf{IDF1 (\%)} & \textbf{HOTA (\%)} & \textbf{IDSW $\downarrow$} & \textbf{RMSE @ 1s (m)} & \textbf{RMSE @ 2s (m)} & \textbf{RMSE @ 3s (m)} & \textbf{Latency (ms)} & \textbf{FPS} \\
\midrule
$M_0$ & YOLOv8n + Standard Kalman SORT (No VIO) & 64.2 & 58.5 & 52.1 & 38 & 0.79 & 1.98 & 3.62 & \textbf{5.1} & \textbf{196.1} \\
$M_1$ & YOLOv8n-P2 + ByteTrack (2D Linear KF, No VIO) & 76.8 & 74.1 & 65.4 & 22 & 0.68 & 1.74 & 3.18 & 5.6 & 178.6 \\
$M_2$ & YOLOv8n-P2 + ByteTrack + VIO Ego-Comp + Lin-KF & 84.5 & 83.2 & 73.8 & 14 & 0.48 & 1.15 & 2.09 & 7.4 & 135.1 \\
\textbf{$M_3$} & \textbf{Proposed: YOLOv8n-P2 + ByteTrack + VIO + 3D CTRA-EKF} & \textbf{89.4} & \textbf{91.2} & \textbf{81.6} & \textbf{8} & \textbf{0.19} & \textbf{0.36} & \textbf{0.61} & 7.8 & 128.2 \\
\bottomrule
\end{tabular}}
\end{table*}

\section{Experimental Results \& Performance Analysis}

\subsection{Small-Drone Detection Evaluation}
Table~\ref{tab:detection} compares our YOLOv8n-P2 detector against five lightweight baselines across all $14,400$ frames. Baseline YOLOv8n achieves only $73.2\%$ Recall and $78.6\%$ $\text{AP}_{50}$ due to feature loss on drones smaller than $16 \times 16\text{ px}$ in Regime $R_1$. Adding the $160 \times 160$ $P_2$ detection head and Wise-IoU regression in our proposed model boosts Recall by $+15.5$ percentage points to $88.7\%$, Precision to $92.3\%$, $\text{F1-Score}$ to $90.5\%$, and $\text{AP}_{50}$ to $91.4\%$ ($\text{mAP}_{50:95} = 58.9\%$), outperforming SOD-YOLO \cite{ref3} by $+3.6\%$ $\text{AP}_{50}$ and YOLO-SR \cite{ref7} by $+3.2\%$ $\text{AP}_{50}$ while requiring only $2.84\text{ M}$ parameters (Fig.~\ref{fig:pr_ablation}a).

\begin{figure*}[!t]
\centering
\includegraphics[width=0.88\textwidth]{fig3_uav_pr_ablation.png}\vspace{-8pt}
\caption{Quantitative evaluation results: (a) Precision-Recall curves comparing YOLOv8n-P2+WIoU against SOD-YOLO \cite{ref3} and baseline YOLOv8n, (b) Multi-object tracking ablation ($M_0\text{--}M_3$) across MOTA, IDF1, and HOTA, and (c) 3D trajectory prediction $\text{RMSE}$ (m) as a function of occlusion duration ($0.5\text{--}3.0\text{ s}$).}
\label{fig:pr_ablation}
\end{figure*}

\subsection{Ablation Study: ByteTrack, VIO Compensation \& CTRA-EKF}
Table~\ref{tab:ablation} and Fig.~\ref{fig:pr_ablation}b--c present the systematic ablation from $M_0$ to $M_3$:
\begin{enumerate}
    \item \textit{Effect of P2 Head \& ByteTrack ($M_0 \to M_1$):} Retaining low-confidence boxes ($0.15 \le s < 0.50$) during partial foliage clutter in Regime $R_2$ improves $\text{IDF1}$ from $58.5\%$ to $74.1\%$ and cuts identity switches ($\text{IDSW}$) from $38$ to $22$.
    \item \textit{Effect of VIO Camera Ego-Motion Compensation ($M_1 \to M_2$):} Under aggressive camera pan/tilt jitter in Regime $R_4$, uncompensated image-plane predictions in $M_1$ suffer frequent bounding-box association failures. Integrating tightly coupled VIO ($\text{ATE}_{\text{RMSE}} = 0.084\text{ m}$, $\text{RPE} = 0.019\text{ m/frame}$) warps tracklet priors accurately via Eq.~(\ref{eq:ego_warp}), elevating $\text{MOTA}$ to $84.5\%$ and $\text{IDF1}$ to $83.2\%$.
    \item \textit{Effect of 3D CTRA-EKF Under Occlusion ($M_2 \to M_3$):} During complete $2.0\text{ s}$ and $3.0\text{ s}$ structural occlusions in Regime $R_3$, constant-velocity linear Kalman filtering ($M_2$) drifts tangentially when the target UAV executes a banked turn ($\text{RMSE} = 1.15\text{ m}$ at $2.0\text{ s}$ and $2.09\text{ m}$ at $3.0\text{ s}$). Replacing the linear model with the 3D CTRA-EKF ($M_3$) reduces prediction $\text{RMSE}$ to $0.19\text{ m}$ at $1.0\text{ s}$, $0.36\text{ m}$ at $2.0\text{ s}$ ($-81.8\%$ vs. $M_0$), and $0.61\text{ m}$ at $3.0\text{ s}$, driving overall $\text{IDF1}$ to $91.2\%$, $\text{HOTA}$ to $81.6\%$, and reducing $\text{IDSW}$ to just $8$ across all 16 sequences (Fig.~\ref{fig:ekf_fps}a).
\end{enumerate}

\begin{figure*}[!t]
\centering
\includegraphics[width=0.86\textwidth]{fig4_uav_ekf_vio_fps.png}\vspace{-8pt}
\caption{State estimation and computational efficiency: (a) Top-down world-frame UAV trajectory reconstruction across $2.2\text{ s}$ and $2.8\text{ s}$ blind occlusion windows comparing Proposed CTRA-EKF+VIO ($\text{RMSE}=0.29\text{ m}$) against uncompensated linear KF ($\text{RMSE}=1.48\text{ m}$), and (b) Stage-wise per-frame latency and real-time FPS across desktop GPU, laptop GPU, and embedded Jetson Orin Nano hardware.}
\label{fig:ekf_fps}
\end{figure*}

\subsection{Computational Efficiency \& Edge Feasibility}
As shown in Fig.~\ref{fig:ekf_fps}b, the complete $M_3$ pipeline requires only $318\text{ MB}$ of GPU VRAM and executes in $7.8\text{ ms/frame}$ ($128.2\text{ FPS}$) on an NVIDIA RTX 3060 desktop GPU ($4.8\text{ ms}$ YOLOv8n-P2 detection, $2.1\text{ ms}$ VIO update, and $0.9\text{ ms}$ ByteTrack + CTRA-EKF), $13.9\text{ ms/frame}$ ($71.9\text{ FPS}$) on a laptop GTX 1650 GPU, and $21.7\text{ ms/frame}$ ($46.1\text{ FPS}$) on an embedded NVIDIA Jetson Orin Nano ($15\text{ W}$ mode)---well within the $33.3\text{ ms}$ ($30\text{ FPS}$) real-time video budget.

\section{Discussion, Limitations \& Future Work}
The ablation results in Table~\ref{tab:ablation} show why isolated improvements in detection or 2D box association fall short during mobile surveillance. Without VIO camera-motion compensation, gimbal pitch jitter alone causes up to $38$ identity switches; likewise, without curvilinear 3D state propagation, blind occlusions longer than $1.5\text{ s}$ cause the tracker to miss banked turns.

\textbf{Practical Limitations:} First, monocular VIO scale observability degrades under pure rotational camera motion lacking translational parallax ($\text{ATE}_{\text{RMSE}}$ increases from $0.084\text{ m}$ to $0.19\text{ m}$ under stationary pan-only sweeps). Second, for occlusions exceeding $4.0\text{ s}$ where the target UAV reverses turn direction mid-occlusion, open-loop CTRA-EKF covariance grows super-linearly. Third, extreme low-light or thermal crossover conditions degrade RGB feature extraction.

\textbf{Future Work:} Future extensions will incorporate RGB-Thermal (RGB-T) cross-modal fusion \cite{ref7}, Interactive Multiple Model (IMM) filtering switching between hover, linear dash, and coordinated turn regimes, and hardware-in-the-loop field flight trials on an onboard RK3588 / Jetson Orin companion computer.

\section{Conclusion}
We have presented and tested a lightweight vision-and-inertial framework for detecting, tracking, and predicting the 3D trajectories of small quadrotors under occlusion and camera ego-motion. By tightly coupling YOLOv8n-P2 small-target detection ($91.4\%$ $\text{AP}_{50}$) with VIO camera ego-motion compensation ($\text{ATE}_{\text{RMSE}} = 0.084\text{ m}$), confidence-stratified ByteTrack association, and a 3D CTRA-EKF predictor in Gazebo Harmonic 8.15, the proposed system achieved $89.4\%$ MOTA, $91.2\%$ IDF1, $81.6\%$ HOTA, and an $81.8\%$ reduction in $2.0\text{ s}$ occlusion prediction error while operating at $128.2\text{ FPS}$ on desktop GPU and $46.1\text{ FPS}$ on embedded edge hardware.

\begin{thebibliography}{20}
\bibitem{ref1} L. Xu, Y. Zhao, Y. Zhai, L. Huang, and C. Ruan, ``Small Object Detection in UAV Images Based on YOLOv8n,'' \textit{Int. J. Comput. Intell. Syst.}, vol. 17, art. 223, 2024, doi: 10.1007/s44196-024-00632-3.
\bibitem{ref2} P. Zhu \textit{et al.}, ``Detection and Tracking Meet Drones Challenge,'' \textit{IEEE Trans. Pattern Anal. Mach. Intell.}, vol. 44, no. 11, pp. 7380--7399, 2022, doi: 10.1109/TPAMI.2021.3119563.
\bibitem{ref3} Y. Li, Q. Li, J. Pan, Y. Zhou, H. Zhu, H. Wei, and C. Liu, ``SOD-YOLO: Small-Object-Detection Algorithm Based on Improved YOLOv8 for UAV Images,'' \textit{Remote Sensing}, vol. 16, no. 16, art. 3057, 2024, doi: 10.3390/rs16163057.
\bibitem{ref4} W. Xu, C. Cui, Y. Ji, X. Li, and S. Li, ``YOLOv8-MPEB Small Target Detection Algorithm Based on UAV Images,'' \textit{Heliyon}, vol. 10, no. 8, art. e29501, 2024, doi: 10.1016/j.heliyon.2024.e29501.
\bibitem{ref5} F. Feng, X. Yang, R. Yang, and F. Zhu, ``Improved YOLOv8 Algorithms for Small Object Detection in Aerial Imagery,'' \textit{J. King Saud Univ. - Comput. Inf. Sci.}, vol. 36, no. 6, art. 102113, 2024, doi: 10.1016/j.jksuci.2024.102113.
\bibitem{ref6} Y. Xiao and N. Di, ``SOD-YOLO: A lightweight small object detection framework,'' \textit{Scientific Reports}, vol. 14, art. 25624, 2024, doi: 10.1038/s41598-024-77513-4.
\bibitem{ref7} S. Liang, X. Feng, M. Xie, Q. Tang, H. Zhu, and G. Li, ``Lightweight YOLO-SR: A Method for Small Object Detection in UAV Aerial Images,'' \textit{Applied Sciences}, vol. 15, no. 24, art. 13063, 2025, doi: 10.3390/app152413063.
\bibitem{ref8} X. Jin, X. Zhu, D. Kang, W. Shen, Y. Zhao, and X. Li, ``GS-YOLO: A lightweight high-accuracy model for small target detection in drone aerial images,'' \textit{PLOS ONE}, vol. 21, no. 6, art. e0350840, 2026, doi: 10.1371/journal.pone.0350840.
\bibitem{ref9} Y. Zhang, P. Sun, Y. Jiang, D. Yu, F. Weng, Z. Yuan, P. Luo, W. Liu, and X. Wang, ``ByteTrack: Multi-Object Tracking by Associating Every Detection Box,'' in \textit{Proc. Eur. Conf. Comput. Vis. (ECCV)}, 2022, pp. 1--21, doi: 10.1007/978-3-031-20047-2\_1.
\bibitem{ref10} J. Cao, J. Pang, X. Weng, R. Khirodkar, and K. Kitani, ``Observation-Centric SORT: Rethinking SORT for Robust Multi-Object Tracking,'' in \textit{Proc. IEEE/CVF Conf. Comput. Vis. Pattern Recognit. (CVPR)}, 2023, pp. 9686--9696.
\newline\vfill\eject
\bibitem{ref11} N. Aharon, R. Orfaig, and B.-Z. Bobrovsky, ``BoT-SORT: Robust Associations Multi-Pedestrian Tracking,'' \textit{arXiv preprint arXiv:2206.14651}, 2022.
\bibitem{ref12} T. Qin, P. Li, and S. Shen, ``VINS-Mono: A Robust and Versatile Monocular Visual-Inertial State Estimator,'' \textit{IEEE Trans. Robotics}, vol. 34, no. 4, pp. 1004--1020, 2018, doi: 10.1109/TRO.2018.2853729.
\bibitem{ref13} G. Welch and G. Bishop, ``An Introduction to the Kalman Filter,'' Dept. Comput. Sci., Univ. North Carolina at Chapel Hill, Tech. Rep. TR 95-041, 1995.
\bibitem{ref14} D. Simon, \textit{Optimal State Estimation: Kalman, $H_\infty$, and Nonlinear Approaches}. Hoboken, NJ, USA: Wiley, 2006.
\bibitem{ref15} G. Jocher, A. Chaurasia, and J. Qiu, ``Ultralytics YOLOv8,'' version 8.0.0, 2023. [Online]. Available: \url{https://github.com/ultralytics/ultralytics}
\bibitem{ref16} J. Luiten \textit{et al.}, ``HOTA: A Higher Order Metric for Evaluating Multi-Object Tracking,'' \textit{Int. J. Comput. Vis.}, vol. 129, pp. 548--578, 2021, doi: 10.1007/s11263-020-01375-2.
\bibitem{ref17} E. Ristani, F. Solera, R. Zou, R. Cucchiara, and C. Tomasi, ``Performance Measures and a Data Set for Multi-Target, Multi-Camera Tracking,'' in \textit{Proc. ECCV Workshops}, 2016, pp. 17--35, doi: 10.1007/978-3-319-48881-3\_2.
\bibitem{ref18} J. Sturm, N. Engelhard, F. Endres, W. Burgard, and D. Cremers, ``A Benchmark for the Evaluation of RGB-D SLAM Systems,'' in \textit{Proc. IEEE/RSJ Int. Conf. Intell. Robots Syst. (IROS)}, 2012, pp. 573--580, doi: 10.1109/IROS.2012.6385773.
\bibitem{ref19} K. Bernardin and R. Stiefelhagen, ``Evaluating Multiple Object Tracking Performance: The CLEAR MOT Metrics,'' \textit{EURASIP J. Image Video Process.}, vol. 2008, art. 246309, 2008, doi: 10.1155/2008/246309.
\bibitem{ref20} Y. Zhang, C. Wang, X. Wang, W. Zeng, and W. Liu, ``FairMOT: On the Fairness of Detection and Re-Identification in Multiple Object Tracking,'' \textit{Int. J. Comput. Vis.}, vol. 129, pp. 3069--3087, 2021.
\end{thebibliography}

\end{document}
